% La Familia: two-page extended abstract, ACM format (acmart, sigconf).
% Build on Overleaf: New Project -> Upload Project (this file), or open the
% "ACM Conference Proceedings Primary Article Template" and replace main.tex with it.
% Compiler: pdfLaTeX. Result must be exactly two pages; see the notes at the end.
\documentclass[sigconf,nonacm]{acmart}

\usepackage{booktabs}

\settopmatter{printacmref=false, printccs=false, printfolios=false}
\pagestyle{plain}

\begin{document}

\title{La Familia: An Offline-First Family Chat App for Android}
\subtitle{Mobile Computing (FNI-421), Phase~1}

\author{Rana Ahmed ElAhmedy}
\affiliation{\institution{Mobile Computing (FNI-421), Team 43}\country{}}
\author{Abdelrhman Salah Ahmed}
\affiliation{\institution{Mobile Computing (FNI-421), Team 43}\country{}}
\author{Alaa Ahmed Yagoub Eissa}
\affiliation{\institution{Mobile Computing (FNI-421), Team 43}\country{}}
\author{Abdelrhman Abdelrahman Hosnie Abdelradi}
\affiliation{\institution{Mobile Computing (FNI-421), Team 43}\country{}}
\author{Youssef Emad}
\affiliation{\institution{Mobile Computing (FNI-421), Team 43}\country{}}
\author{Muhammed Kamel}
\affiliation{\institution{Mobile Computing (FNI-421), Team 43}\country{}}

\renewcommand{\shortauthors}{ElAhmedy et al.}

\begin{abstract}
Messaging apps assume a stable connection, yet many families live with slow or
interrupted mobile data. La Familia is an Android chat app in which a conversation
never depends on the network being available at the moment a message is written.
It is built with Kotlin and Jetpack Compose; a Room database is the single source of
truth, WorkManager delivers queued messages with exponential back-off, and Firebase
Firestore synchronises devices in real time. The app supports accounts, user search,
one-to-one chats with edit, delete and retry, delivered and read receipts,
per-user pin, mute and archive, English and Arabic with right-to-left layout, light and
dark themes saved with DataStore, and sharing through implicit Intents. Business rules
are pure Kotlin functions covered by 51 unit tests, and the project passes Android lint
with no warnings.
\end{abstract}

\keywords{Android, Jetpack Compose, Room, WorkManager, Firebase Firestore, DataStore,
offline-first, Intents, localisation}

\maketitle

\section{Problem}
Popular messengers treat the network as always present: when the connection drops, the
composer still accepts text but the message silently waits, fails, or is lost when the
app is closed. Families who share a single data plan or live in areas with weak coverage
need the opposite behaviour. A message written offline should be saved at once, shown
with an honest status, and delivered automatically later, and the history of every
conversation should be readable without a connection. Our goal for Phase~1 was a
complete, small messenger that demonstrates this offline-first design together with the
course topics: a Compose user interface, Intents, persistence with full CRUD, and clean,
tested Kotlin.

\section{Solution and Architecture}
La Familia is a single-activity app with unidirectional data flow. Each screen is a
composable driven by a \texttt{ViewModel} that exposes one immutable UI-state object
through a Kotlin \texttt{StateFlow}; user actions are plain function calls back into the
ViewModel. Screens only know a \texttt{ChatRepository} interface, which let six team
members build their parts in parallel against one contract.

\textbf{Room as the single source of truth.} The production repository reads
conversations and messages only from Room. DAO queries return \texttt{Flow}s, so the
user interface redraws itself whenever a row changes, whether the change came from the
user or from the server.

\textbf{Outbox with WorkManager.} Sending a message inserts it into Room with the status
\textsc{queued} and enqueues a unique WorkManager job with a network constraint and
exponential back-off. The worker writes the message to Firestore in a batch with the
conversation's last-message preview. Because the document id is fixed, a retry can never
create a duplicate. After repeated failures the message becomes \textsc{failed} and the
user can tap it to retry.

\textbf{Firestore to Room.} A snapshot listener on the user's conversations, and one on
each conversation's messages, copies additions, edits and deletions into Room. The
receiving phone writes a \texttt{deliveredAt} timestamp when a message arrives and a
\texttt{readAt} timestamp when the chat is opened; the sender's phone turns them into
\textsc{delivered} and \textsc{read}. A pure rule only lets a status move forward, so a
late update can never turn ``read'' back into ``delivered''.

\textbf{Security.} Firestore rules allow a signed-in user to read public profiles and to
change only their own. Conversations and messages are visible only to members, only the
sender may edit or delete a message, and the other member may add nothing except the two
receipt fields.

\section{Features}
Table~\ref{tab:features} maps the features to the course topics they demonstrate.

\begin{table}[h]
  \caption{Features and the topics they cover}
  \label{tab:features}
  \small
  \begin{tabular}{@{}p{0.30\linewidth}p{0.64\linewidth}@{}}
    \toprule
    Topic & Implementation in La Familia \\
    \midrule
    Compose UI & Six screens in Material~3 (login, register, chat list, chat thread,
                 profile, settings) plus a new-chat dialog; navigation with an auth gate \\
    Intents & Invite through the share sheet (\texttt{ACTION\_SEND} with a chooser) and
              feedback e-mail (\texttt{ACTION\_SENDTO}, \texttt{mailto:}) \\
    Room CRUD & Create, read, edit and delete messages; create, update and remove
                conversations; schema migration from version~1 to~2 \\
    DataStore & Theme, language and notification settings, with reset \\
    Background work & WorkManager outbox with constraints and back-off \\
    Firebase & E-mail/password authentication and real-time Firestore sync \\
    Localisation & English and Arabic with automatic right-to-left mirroring \\
    \bottomrule
  \end{tabular}
\end{table}

\textbf{Accounts.} Registration and login use Firebase Authentication with input
validation and a password-strength hint. A signed-in user skips the login screen, and
each user has a public profile (name and status) that others can search by name.

\textbf{Conversations.} The list shows the last message, its time and an unread badge,
and can be searched. A long press opens a menu to pin, mute, archive or remove a chat.
These choices are stored per user, so two members of the same conversation can organise
it differently. Pinned chats stay on top and archived chats move to their own section.

\textbf{Messages.} Bubbles show the time and a status with ticks (one tick when the
server has the message, two when it reached the other phone, two in the accent colour
when it was read). Users can edit or delete their own messages.

\textbf{Personalisation.} Language and theme can be changed at runtime without
restarting the app. Choosing Arabic mirrors every screen to right-to-left.

\section{Implementation and Results}
The app targets Android~8.0 (API~26) and newer and is written entirely in Kotlin, with
Jetpack Compose, Navigation Compose, Lifecycle ViewModels, Room with KSP, WorkManager,
Preferences DataStore and the Firebase BoM.

Rules that do not need Android, such as input validation, password strength, who may
edit a message, time formatting, conversation filtering and ordering, entity mapping and
status transitions, are pure Kotlin functions. They are covered by 51 JVM unit tests,
together with an in-memory fake of the repository used to test create, read, update and
delete without a device. All tests pass, Android lint reports no issues, and the build
has no compiler warnings. The code uses \texttt{val} by default, no non-null assertions
(\texttt{!!}), and Kotlin idioms such as scope functions, sequences and KTX helpers.

The work was split into six parts, one per member (foundation and accounts; conversation
list; chat thread; Room, Firestore and WorkManager; settings, theme and localisation;
Intents, tests and documentation). Each part was delivered as its own pull request
against the shared repository contract and reviewed before merging. Because every read
comes from Room, conversations stay readable in airplane mode; a message written offline
shows ``waiting'' and is delivered automatically when the connection returns, and
settings survive restarts.

\section{Conclusion}
La Familia shows that an offline-first design is practical for a small team: treating
the local database as the source of truth and the network as a synchronisation channel
makes the app fast, honest about message status, and usable with poor connectivity.
Future work includes group chats, push notifications, typing indicators, an offline
banner and layouts tuned for tablets.

% ---------------------------------------------------------------------------
% Notes for the team (not printed):
% * The course asks for EXACTLY two pages. Compile and check the page count:
%   - more than two pages: shorten "Features" or the table first;
%   - less than two pages: add a figure with two app screenshots, e.g.
%       \begin{figure}[h]\centering
%         \includegraphics[width=0.48\linewidth]{chatlist.png}\hfill
%         \includegraphics[width=0.48\linewidth]{chat.png}
%         \caption{Chat list and chat thread (light theme).}\end{figure}
%     (upload the PNGs to Overleaf next to this file).
% * Two members must peer-review the text before the PDF is added to the repo.
% * Before submitting, confirm on a device what the text claims: airplane mode keeps chats
%   readable, an offline message shows "waiting" then sends, settings survive a restart,
%   and receipts / pin / mute / archive work (PRs #95 and #96 must be merged).
\end{document}
